GLP-1 receptor agonists are an important class of medications used to treat diabetes, cardiovascular disease, and obesity \citep{Drucker2018}. In randomized trials, they generate weight loss, improve glycemic control, prevent transitions from pre-diabetes to diabetes, and reduce major adverse cardiac events and mortality \citep{Drucker2018, Marso2016a, Wilding2021, Jastreboff2022, lincoff2023semaglutide, le2017liraglutide, kahn2024semaglutide, jastreboff2025tirzepatide}. Different dosages of the same GLP-1 molecule are marketed under different brand names for different FDA-approved uses---e.g., semaglutide is sold as Ozempic for diabetes and Wegovy for obesity. The purpose of the drug matters for coverage in public insurance programs. Nearly all FDA-approved drugs---including GLP-1 medications for diabetes and cardiovascular disease---are covered by Medicare Part D and Medicaid. Under federal law, however, Medicare Part D plans \emph{must exclude} and state Medicaid plans \emph{may exclude} FDA-approved drugs if they are "used for anorexia, weight loss, or weight gain" \citep{SSA1860D2, SSA1927}. Policy makers have pushed to make the drugs more accessible by negotiating lower prices for GLP-1 medications for public insurance programs and consumers \citep{cms2025balance_page, WhiteHouseTrumpRx2025}. Federal legislative proposals have sought to expand coverage of obesity medications under Medicare Part D and Medicaid but have not become law \citep{TROA2025}. In 2026, a Medicare Part D final 
rule declined to reinterpret the statutory exclusion administratively \citep{CMSFinalRule2026}, though CMS has since launched demonstration programs to expand access \citep{cms2025balance_page}.

Despite their clinical benefits, GLP-1 medications are unlikely to pay for themselves over the short to medium run. Recent research suggests that the drugs are improving health in real-world settings \citep{zhang2026glp1, lu2025changes}, but treating patients with GLP-1 medications increases other health care spending over 1--5 year horizons \citep{wing2026glp1, wennberg2025glp1}. This makes the question of who pays---and how much---central to the policy debate over expanding access.

Because coverage is discretionary, there is substantial variation in whether state Medicaid programs pay for obesity-indicated GLP-1 medications. Previous research examines cross-sectional variation in state Medicaid coverage for GLP-1 medications, but does not measure the effects of coverage changes on utilization and costs \citep{liu2024state, klebanoff2025medicaid}. Since 2022, seventeen states have offered Medicaid coverage for at least one obesity-indicated GLP-1 medication. However, four of these states recently rescinded coverage, citing budget pressures and unsustainable pharmacy costs. There are good reasons to expect the fiscal burden to be large. Around 40 percent of adult Medicaid enrollees meet clinical criteria for obesity, and 70 percent meet criteria for overweight or obesity \citep{mylona2020association}. A large literature shows that insurance coverage increases health care utilization \citep{manning1987health,finkelstein2012oregon}, suggesting that extending Medicaid coverage for GLP-1 medications to this population could generate substantial new costs for state budgets. Such large changes in Medicaid spending may strain or distort federal and state budget priorities in ways not directly related to health or health care \citep{baicker2005fiscal, gruber2017federal, clemens2018implications}. The Medicaid Drug Rebate Program (MDRP) requires manufacturers to pay substantial rebates to states, which could significantly reduce the net cost of coverage. However, the effects of these rebates on state government costs are not well understood, and some research suggests that substantial Medicaid purchases can drive up the drug prices paid by non-Medicaid consumers as manufacturers set prices and design products strategically \citep{duggan2006distortionary}.

We make three contributions. First, we compile data on the timing and details of Medicaid coverage for obesity-indicated GLP-1 medications across states. Second, using a stacked difference-in-differences estimator that identifies an enrollment-weighted average treatment effect for early adopter states, we show that coverage increased Medicaid prescriptions by \rxObSevenNineEmw{} per 100 enrollee-months nine quarters after adoption, and did not affect GLP-1 prescribing for diabetes or cardiovascular indications. This suggests that coverage did not induce substitution away from diabetes formulations, which could have been used off-label to treat obesity. We also find no evidence that Medicaid coverage displaced private out-of-pocket purchases from GLP-1 compounding platforms, which we measure using credit and debit card transactions data. This suggests that Medicaid coverage generated new GLP-1 utilization rather than crowding out existing private purchases. Third, we estimate the net fiscal cost of coverage to state budgets, accounting for manufacturer rebates under the MDRP. Nine quarters after adoption, coverage increased net-of-rebate Medicaid expenditures by \$\netSpObSevenNineEmw{} per 100 enrollee-months, compared to gross reimbursements of \$\spObSevenNineEmw{}. The results suggest that MDRP rebates for obesity-indicated GLP-1 drugs reduce program costs by about \spReductionObSevenNineEmw{} percent and that manufacturer best prices and drug specific inflation are not key drivers of the rebates for these drugs. 

Our estimates help put decisions about Medicaid coverage of GLP-1 drugs into perspective. California, for example, started covering obesity-indicated GLP-1 medications in 2023 and rescinded coverage in 2026. There were \enrollCA{} people enrolled in California's Medicaid program in September 2025. Our nine-quarters out utilization effect estimates imply that by covering the drugs, California's Medicaid program was paying for about \usersCAetNine{} people to receive obesity-indicated GLP-1 medications each month, with an annual net-of-rebate cost of \$\netCostCAetNine{}. Thus, the savings from rescinding coverage are not small. Aggregating across all seventeen coverage states suggests that Medicaid coverage for obesity-indicated GLP-1s costs about \$\netCostCoveringEtNine{} per year. Extending coverage to the thirty-three states that do not offer coverage would generate another \$\netCostNoncoveringEtNine{} in annual net spending.